\section{Nuisance parameters of the electronic rate}\label{sec:10.10}

A dark-matter particle of \SI{3}{MeV} with a heavy mediator and $\sigmae = \SI{e-37}{cm^2}$ makes $5.2\times10^4$ electron--hole pairs per kilogram-year in silicon if the electrons are independent, $2.5\times10^4$ if they screen in the random-phase approximation, and $1.85\times10^4$ once local fields are included. All three numbers come from the same published tables, counted from the low-temperature gap of \SI{1.17}{eV}, and they differ only in the many-body approximation. The last one is the best computed so far, but it is not the rate. The electron--hole attraction raises it, the conduction bands above the edge are uncertain at the level of a tenth of an eV, and two published tabulations of the same independent-particle response disagree by a third at this mass. \Cref{sec:8.8} turned each truncation into a nuisance parameter $\vartheta_k$, with the computed or estimated step to the next order as its template.

\subsection{Many-body spread}\label{sec:10.10-manybody}

The hierarchy of approximations of \cref{sec:4.6} is not an expansion in a small parameter, and its steps do not shrink. At \SI{3}{MeV} with a heavy mediator, random-phase screening lowers the independent-particle rate by \SI{52}{\percent}, and local fields lower what remains by a further \SI{26}{\percent}; for a light mediator the steps are \SI{76}{\percent} and \SI{18}{\percent}. The ratio of the independent-particle rate to the screened rate with local fields is 2.8 for a heavy mediator and 5.1 for a light one, between the values 2.3 and 9 that \cref{sec:10.2} found at 10 and \SI{1}{MeV}. A spread of this kind is not an uncertainty of the best approximation. Once an approximation is computed it becomes the central value, and the nuisance that remains is the step after it (\cref{sec:8.8}). For silicon that step is the electron--hole attraction, the excitonic (vertex) correction of the Bethe--Salpeter equation, and no finite-momentum calculation of it exists for silicon (\cref{sec:10.5}).

Its size can be estimated from the model of \cref{fig:10-xq}. At $|\vect{q}| = \SI{0.91}{keV}$ the attraction raises the screened loss function by 27 to \SI{29}{\percent} from the onset to \SI{3.5}{eV}, by \SI{49}{\percent} at \SI{4}{eV}, and by about \SI{30}{\percent} above, where part of the rise is an artifact of four conduction bands. A particle of \SI{3}{MeV} deposits at most $\tfrac12\mchi v_{\mathrm{max}}^2 = \SI{10.5}{eV}$, and the mean energy of its spectrum is \SI{4.7}{eV}. We take $\delta_{\mathcal{X}} = +0.28$ over that window, an estimate from a model at one momentum transfer, with the coupling block of \cref{eq:10.5-coupling} adding about $\pm0.10$ to it. The template has a definite sign: at $\vartheta_{\mathcal{X}} = 1$ the rate rises. By the $f$-sum identity \cref{eq:8.8-fsum} a many-body template is a tilt. Inside the \SI{10.5}{eV} window it raises the spectrum nearly uniformly and so acts like a normalization; the weight it moves comes from energies above the window, which the particle cannot reach. The same template, read at \SI{10}{MeV}, reaches \SI{35}{eV} and samples part of the compensating loss, so one $\vartheta_{\mathcal{X}}$ changes the rates at 3 and \SI{10}{MeV} by different fractions. The attraction also adds new support. The bound lines below the pair continuum carry the weight of the free continuum within $8.0\,E_x^b = \SI{0.12}{eV}$ of the edge (\cref{sec:10.5-lines}), which in the tables holds $4\times10^{-4}$ of the heavy-mediator rate at \SI{3}{MeV} and $1.5\times10^{-3}$ of the light-mediator rate. The tables resolve energies to \SI{0.2}{eV}, so these shares are orders of magnitude, and they matter only for thresholds at the gap.

A second nuisance sits inside the approximation already computed. At momentum transfers of two to three times the shortest reciprocal-lattice vector the local fields lower the per-$q$ rate of the tables by 1.67 at \SI{8}{keV} and 1.83 at \SI{10}{keV} (\cref{sec:10.2-corrections}). That size rests on a spectrum that holds only \SI{17}{\percent} of the $f$-sum below \SI{89}{eV} at \SI{10}{keV}. The model of \cref{fig:10-fcrys-limit}, built from a local pseudopotential with no input in common with the tables but the crystal structure, gives the same suppression. At \SI{5}{eV} and between 2 and $2.8\,q_0$ it finds the screened structure factor with local fields at 0.54 to 0.61 of the independent-particle one, with the head screening already within \SI{10}{\percent} of one above $2.2\,q_0$. The tables at the same energy and momenta give 0.40 to 0.48, with the head screening at 0.77 to 0.97 of the independent-particle value. Two calculations that share neither potential nor basis therefore agree that local fields hold the response near half of its unscreened value there, to within a quarter of each other. Both, however, lack the core-region oscillations of the valence states and the high-energy tail of the spectrum, and an all-electron calculation with enough bands is needed to fix the size (\cref{app:E.2,app:E.6}). The template is the change of the rate if the suppression above $2q_0 = \SI{7.46}{keV}$ were absent, $\vartheta = 1$ restoring the random-phase value there. It is $+0.067$ at \SI{3}{MeV} for a heavy mediator, where only \SI{7.5}{\percent} of the rate comes from those momenta, $+0.009$ for a light one, and $+0.24$ at \SI{10}{MeV}, where they hold \SI{29}{\percent} of the rate.

A third source is not a truncation of the label but of the numerics. The crystal form factor of \cref{sec:10.1} and the independent-particle head of the tables of \cref{sec:10.2} are the same quantity in the same approximation, computed with different codes, pseudopotentials and grids. With the same halo, each from its own onset, they give $1.60\times10^5$ and $1.36\times10^5$ events per kilogram-year at \SI{10}{MeV} for a heavy mediator, a ratio of 1.18, and 1.29 for a light one. At \SI{3}{MeV} the ratio is 1.34 for a heavy mediator and 1.55 for a light one. The difference sits at small $q$, where the lighter particle collects its rate, and the $f$-sum rule locates it. Between $0.3\,q_0$ and $0.7\,q_0$ the first moment of the form-factor table is 1.1 to 1.2 times the valence value (\cref{sec:10.1}), and that of the independent-particle head of the tables is 0.93 at every such momentum, partly because those tables end at \SI{89}{eV}. Neither table conserves the $f$-sum at small $q$, and the two violations lie on either side of it. The spread is measured with independent particles, but the central value is screened, and screening is strongest exactly at the small $q$ where the tables disagree. There a normalization error of $\varepsilon_2$ is partly screened out: if $\varepsilon_2$ and, by Kramers--Kronig, $\varepsilon_1 - 1$ are too large by $1 + \delta$ where $\varepsilon_1 \gg 1$, then $L \simeq \varepsilon_2/\varepsilon_1^2$ moves by about $1/(1 + \delta)$, opposite in sign and no larger. The independent-particle ratio is therefore an upper bound on the spread of the screened rate, and we carry it as such. \Cref{sec:8.8} named this case: for a truncation that breaks the $f$-sum the violation measures the error, and the spread between the calculations is the nuisance until a converged calculation passes the test (\cref{app:E.7}).

The Bethe--Salpeter step carries one more pitfall, an evaluation error rather than a truncation. A Bethe--Salpeter calculation that keeps only the resonant poles has the right $\varepsilon_2$ and the wrong $\varepsilon_1$, and it overstates the loss function by 2.12 at \SI{1.39}{keV} and \SI{1.3}{eV}, falling to within \SI{15}{\percent} of one above \SI{6.7}{keV} (\cref{eq:10.5-res}). Both time orderings remove it. It is listed below because a rate that inherits it is wrong by a factor that no nuisance of the correct calculation covers.

\subsection{Gap and threshold nuisances}\label{sec:10.10-gap}

The tables carry a scissor shift that places the conduction bands at the measured gap of silicon at zero temperature, \SI{1.17}{eV}~\cite[Sect.~III\,B]{knapen2021dark}, the gap a detector at millikelvin temperatures has, with the zero-point renormalization included in the measurement (\cref{sec:3.4}). The rate is counted from that gap, and the threshold mass is $m_{\mathrm{min}} = 2\Eg/v_{\mathrm{max}}^2 = \SI{0.333}{MeV}$. The \SI{1.11}{eV} that often accompanies these tables in analysis codes is a different number, the gap of the step model of pair creation with \SI{3.6}{eV} per pair (\cref{sec:1.5}), not the gap of the response. The energy nodes of the tables lie \SI{0.2}{eV} apart and do not resolve the edge itself.

What remains depends on what is measured. The edge itself is known to a few meV, and a shift of \SI{5}{meV} changes the \SI{3}{MeV} rate by \SI{0.4}{\percent}. A scissor, however, moves every conduction band by the same amount, and quasiparticle corrections in silicon differ from state to state by about \SI{0.1}{eV} (\cref{sec:3.4}). The bands above the edge are therefore uncertain at that level even when the edge is not. A shift of the whole spectrum by \SI{0.1}{eV} bounds this effect from above: it changes the \SI{3}{MeV} rate by $-7.5$ and $+8.0$ percent for a heavy mediator, by $\mp10$ percent for a light one, and the \SI{10}{MeV} heavy-mediator rate by \SI{4}{\percent}. The effect grows toward the threshold mass, where the spectrum samples only the first eV above the edge. For a target whose gap has not been measured the edge itself carries the \SI{0.1}{eV}, and $\delta\ln m_{\mathrm{min}} = \delta\ln\Eg$ turns it into \SI{8.5}{\percent} in the threshold mass of silicon (\cref{we:3-mmin-gap}). Near that mass the halo supply changes by orders of magnitude over a factor of two in $\mchi$, and the rate there depends on the gap more than on any other entry of the table below. The exciton binding of \SI{14.7}{meV} is a threshold nuisance of the same kind, which lowers the first line to \SI{1.155}{eV} (\cref{we:10-xline}), and \cref{tab:8-nuisances} carries it.

\subsection{Nuisance table of the electronic sector}\label{sec:10.10-table}

\Cref{tab:10-nuisances} lists the nuisances of the electronic sector. A next order computed in a model gives a template whose sign and shape are known and whose value may change once the real calculation is done. An estimate from a small parameter gives the size to a factor of order one. A spread between published calculations at the same order stands in for a template where the error lies outside the truncation label. An inequality gives a range without a central value, and for some items only the omission can be named.
\begin{table}[tbp]
  \caption{Electronic nuisances. Type in the classes of \cref{sec:8.8}: normalization, tilt (in $\omega$ or in $q$), threshold, new support, and class mix (weight moved between the zero-phonon and phonon-dressed sectors without changing the total). Basis of the size: ``model'' (next order computed in a model), ``tables'' (evaluated from the published tables), ``est.''\ (small-parameter estimate), ``spread'' (published calculations at the same order), ``bound'' (only an inequality is known), ``named'' (no size). Sizes are for silicon unless germanium is named, with rates at \SI{3}{MeV} for a heavy mediator. The last column gives the section where the open calculation is stated.}
  \label{tab:10-nuisances}
  \centering\small\setlength{\tabcolsep}{3pt}
  \begin{tabular}{@{}>{\raggedright\arraybackslash}p{0.26\linewidth}>{\raggedright\arraybackslash}p{0.13\linewidth}>{\raggedright\arraybackslash}p{0.34\linewidth}>{\raggedright\arraybackslash}p{0.09\linewidth}>{\raggedright\arraybackslash}p{0.12\linewidth}@{}}
    \toprule
    Truncation & Type & Size & Basis & Section \\
    \midrule
    Electron--hole attraction ($\mathcal{X}$: RPA + LF $\to$ BSE) & tilt; new support & $L$ up by 27--\SI{29}{\percent} near the edge at \SI{0.91}{keV}; rate $\delta_{\mathcal{X}} \approx +0.28$; lines $4\times10^{-4}$ of the rate & model & \cref{sec:10.5} \\
    Coupling block dropped (BSE in Tamm--Dancoff form) & tilt & about \SI{10}{\percent} of the proper $\varepsilon_2$ & est. & \cref{sec:10.5} \\
    Resonant-only real part (evaluation error) & tilt in $q$ & $L$ $\times2.12$ at \SI{1.39}{keV}, \SI{1.3}{eV}; within \SI{15}{\percent} above \SI{6.7}{keV} & tables & \cref{sec:10.5} \\
    Local fields at $q \gtrsim 2q_0$ (bands, pseudo-states) & tilt in $q$ & per-$q$ factor 1.67 (\SI{8}{keV}), 1.83 (\SI{10}{keV}); rate $+0.066$ if absent above $2q_0$ ($+0.24$ at \SI{10}{MeV}) & spread & \cref{sec:10.2} \\
    Independent-particle tabulation (code, grid, pseudopotential) & normali\-zation at small $q$ & IP rates differ by 1.34 (1.18 at \SI{10}{MeV}), an upper bound for the screened rate; $f$-sum 1.1--1.2 against 0.93 at $0.3$--$0.7\,q_0$ & spread & \cref{sec:10.1} \\
    Quasiparticle gap (rigid scissor) & threshold & edge measured, few meV; bands \SI{0.1}{eV}: rate $-7.5$/$+8.0\,\%$; $m_{\mathrm{min}}$ \SI{8.5}{\percent} without a measured gap & spread & \cref{sec:3.4} \\
    Plasmon decay into pairs with phonons ($n_g = 2$) & class mix & branching $10^{-3}$ to $10^{-2}$ & est. & \cref{sec:10.3} \\
    Local fields of the transverse current & normali\-zation & $-4$ to $-15\,\%$ (model); $-10$ to $-30\,\%$ (est.) & model & \cref{sec:10.4} \\
    Density--current cross term of a velocity coupling & normali\-zation & rate at least 16 (IP), about 29 (screened) times the static-electron formula at \SI{10}{MeV} & bound & \cref{sec:10.4} \\
    Triplet response beyond independent pairs & tilt & weight $\times2.2$ below \SI{3}{eV} (model); exchange splitting $\sim10^{-3}$ of the binding & model & \cref{sec:10.5} \\
    Final states of the germanium $3d$ shell & normali\-zation & $3d$/valence 0.07 (plane waves) against 0.69 (Fermi factor) at \SI{30}{MeV} & bound & \cref{sec:10.6} \\
    Screening of core transitions & normali\-zation & macroscopic \SI{2}{\percent} at \SI{12}{keV} (Ge); local fields not computed & est. & \cref{sec:10.6} \\
    Relativistic corrections (Pauli order) & normali\-zation & about $10^{-3}$, at most \SI{2}{\percent} (Ge $3d$); next order $3\times10^{-6}$ & est. & \cref{sec:10.6} \\
    Spin--orbit mixing at the germanium edge & tilt & up to \SI{10}{\percent} in amplitude & est. & \cref{sec:10.6} \\
    Terms of order $p_\phi^2$ in scalar absorption & new support near $\omega_p$ & density row equal to the kinetic rate at 16--\SI{18}{eV} (model) & named & \cref{sec:10.7} \\
    Room-temperature optical data near the edge ($n_T$) & threshold & more than $\times2$ within about \SI{20}{meV} of the edge & est. & \cref{sec:13.5} \\
    Bound shells in dark Compton scattering & tilt & K shell of Si at 45--\SI{100}{keV}, of Ge up to \SI{118}{keV}; L shells below \SI{20}{keV} & named & \cref{sec:10.8} \\
    Axial pair absorption with spin--orbit coupling & new support & of order $v^4$ without it & named & \cref{sec:10.8} \\
    Average over the directions of $\vect{q}$ & tilt in time & first anisotropic term $\ell = 4$ & named & \cref{sec:11.7} \\
    \bottomrule
  \end{tabular}
\end{table}

Two of these nuisances change the silicon rate at \SI{3}{MeV} by more than ten percent: the electron--hole attraction, by about a quarter, and the spread between tabulations of the independent-particle response, by a third for a heavy mediator and by half for a light one. At \SI{10}{MeV} the local fields at large momentum transfer join them, with a quarter, and the tabulation spread falls to a fifth. The gap stays below ten percent above \SI{3}{MeV} but grows without limit toward the threshold mass. Outside the density rows the largest items are not nuisances in the sense of a computed step. The cross term of a velocity coupling and the final states of the germanium $3d$ shell are known only between bounds, a factor of ten apart for germanium near \SI{30}{MeV}. A limit set on such a coupling has no central value yet, and the bound itself is what it can use (\cref{sec:10.4}).

One rule governs how the rows enter a likelihood. A truncation that two rates share acts on both with one $\vartheta_k$ (\cref{sec:8.8}). The attraction raises the rate at every mass and for both mediators at once, by different fractions, and a combination of masses or experiments that treated it as independent nuisances would let the data play one against the other. The inelastic, exothermic, boosted and fermionic rates of \cref{sec:10.9} read the same loss function and transverse current at shifted kinematics, and they carry the same nuisances with the same parameters; they add no rows of their own. Where one rate uses the tables and another the measured optical constants, as the scattering and absorption rates of \cref{sec:10.2,sec:10.7} do, the two share the many-body truncation only through what the measurement contains, and the absorption rate carries the temperature row instead.

\begin{workedexample}{The nuisance band on the silicon rate at \SI{3}{MeV}}\label{we:10-band3}
Take $\mchi = \SI{3}{MeV}$, a heavy mediator, $\sigmae = \SI{e-37}{cm^2}$ and the standard halo of \cref{sec:1.2}, with $\rho_\chi = \SI{0.3}{GeV.cm^{-3}}$. The kinematic window ends at $\omega = \SI{10.5}{eV}$ and $q = 2\mchi v_{\mathrm{max}} = \SI{15.9}{keV}$. The central value is the loss-function rate \cref{eq:10.2-rate} with the random-phase loss function with local fields of the tables of \cref{sec:10.2}, counted from the low-temperature gap of \SI{1.17}{eV}:
\begin{equation*}
  \mathcal{R}^{\mathfrak{T}} = \SI{1.85e4}{kg^{-1}.yr^{-1}} ,
\end{equation*}
about 51 events per kilogram and day. Its spectrum has a mean energy of \SI{4.7}{eV}, and \SI{10}{\percent} of it lies below \SI{3}{eV}.

Four nuisances act on this number with sizes in hand. In the power form \cref{eq:8.8-M}, $\mathcal{R}(\vartheta) = \mathcal{R}^{\mathfrak{T}}\prod_k(1 + \delta_k)^{\vartheta_k}$, each enters through $\ln(1 + \delta_k)$:
\begin{equation*}
  \begin{aligned}
  &\text{attraction, } \delta_{\mathcal{X}} = +0.28: & \ln(1 + \delta) &= +0.247, \\
  &\text{tabulation spread (bound), } +0.335: & &+0.289, \\
  &\text{band energies, } \SI{0.1}{eV}: & &-0.078, \\
  &\text{local fields above } 2q_0,\ +0.066: & &+0.064.
  \end{aligned}
\end{equation*}
The four are separate truncations with independent standard-normal $\vartheta_k$. To first order the logarithm of the rate then has the Gaussian width of \cref{eq:C.14-quadrature}, the four logarithms added in quadrature,
\begin{equation*}
  \sigma_{\ln\mathcal{R}} = \bigl(0.247^2 + 0.289^2 + 0.078^2 + 0.064^2\bigr)^{1/2} = 0.393 ,
\end{equation*}
and the rate lies between $\mathcal{R}^{\mathfrak{T}}/\ee^{0.393}$ and $\mathcal{R}^{\mathfrak{T}}\ee^{0.393}$ at one standard deviation:
\begin{equation*}
  \SI{1.25e4}{kg^{-1}.yr^{-1}} \le \mathcal{R} \le \SI{2.75e4}{kg^{-1}.yr^{-1}} ,
\end{equation*}
a factor of 1.48 either way. The attraction and the tabulation spread carry \SI{93}{\percent} of the variance. Taking the attraction at the upper value of the model, $+0.49$, widens the factor to 1.65; the coupling block, $\pm0.10$ on $\delta_{\mathcal{X}}$, changes $\ln(1 + \delta_{\mathcal{X}})$ by 0.075.

The band is not symmetric about the best guess. The attraction has a known sign, and $\vartheta_{\mathcal{X}} = 1$ moves the central value to $2.37\times10^4$; the band of the power form describes the uncertainty of where the next approximation lands, not a symmetric spread around the present one. The tabulation spread has no known sign, since the two tables miss the $f$-sum on opposite sides. For a light mediator the same steps give a central value of $6.0\times10^3$ and a factor of 1.67, with the tabulation spread at $+0.55$ and the gap at $-0.10$, and the local fields at large $q$ below one percent.

The band covers the truncations that the density response of silicon carries at this mass. It does not cover the halo, which the likelihood of \cref{eq:8.8-like} holds fixed, nor the detector, whose nuisances \cref{sec:21.4} adds. It is also wider than the step between the two best computed approximations, $2.49\times10^4$ to $1.85\times10^4$. That step is not an uncertainty, since it has been taken; the band covers the next step and the spread between calculations in the same approximation.
\end{workedexample}

\subsection{Open problems}\label{sec:10.10-open}

A finite-momentum Bethe--Salpeter calculation for silicon and germanium, with the coupling block kept and the bound lines added from computed masses and vertices, fixes the attraction, the coupling block, the triplet channel and the spin--orbit mixing at the germanium edge (\cref{sec:10.5,sec:10.6}; \cref{app:E.15}). An all-electron random-phase calculation with local fields, with enough bands to hold the $f$-sum at \SI{10}{keV} and with the current and kinetic rows carried as external rows of the Dyson equation, fixes the local fields at large $q$, the screening of core transitions, the transverse-current fields, the in-medium ratio of velocity couplings and the order-$p_\phi^2$ terms of scalar absorption (\cref{sec:10.2,sec:10.4,sec:10.7}; \cref{app:E.6,app:E.16}). Until these rows are computed, \cref{eq:10.4-band} is all a limit on a velocity coupling can use. Crystal conduction states for the germanium $3d$ shell, and a four-component response for the bound on relativistic terms, settle the core rows (\cref{sec:10.6}). A second-order electron--phonon calculation at the plasma frequency, and shell Compton profiles, settle the plasmon branching and the bound-shell correction to dark Compton scattering (\cref{sec:10.3,sec:10.8}; \cref{app:E.26}). The spread between tabulations needs no new physics, only one converged calculation that passes the $f$-sum test at small $q$.

The quadratic response, two published rates and the self-trapped excitons of NaI are open items that are not nuisances of the rates computed here. The interference between a contact amplitude and a second-order one needs the quadratic response of the crystal, which no reduction to two-point functions supplies (\cref{sec:10.8-open}). The free-electron vertex of fermionic absorption on a bound shell underestimates the rate by at least $1/r$, and the density row of cosmic-ray-boosted dark matter without its retardation factor overstates the rate at relativistic speeds (\cref{sec:10.9-open}). And the self-trapped excitons of NaI, where the Wannier picture fails, are open in every quantity that the rate needs (\cref{sec:10.5}).

\section*{Summary}

An electron--hole pair in silicon costs at least the gap, and the rate of making pairs is a contraction of the master formula with the response of the electrons (\cref{sec:10.1}). For independent electrons the structure factor is a sum over valence--conduction pairs with plane-wave vertices at the unfolded momentum transfer, and its angle average, rescaled, is the crystal form factor on which most published limits rest. Its normalization is where tabulations go wrong, and the $f$-sum rule tests it. From the published silicon form factor a \SI{10}{MeV} particle with a heavy mediator and $\sigmae = \SI{e-37}{cm^2}$ gives $1.60\times10^5$ events per kilogram-year, \SI{27}{\percent} of them from $q > 2q_0$, where pseudo-wavefunctions lose accuracy. The excess of its $f$-sum at small $q$ is only partly explained by the scissor shift.

Screening enters through one function, the loss function, and the loss-function rate is exact for a coupling to the total charge and for a coupling to the electron density at clamped nuclei or above the phonon band (\cref{sec:10.2}). In the published silicon tables the \SI{10}{MeV} rate is $1.36\times10^5$ for independent particles (two independent-particle tabulations differ by \SI{18}{\percent}, \cref{sec:10.10-manybody}), $8.6\times10^4$ with random-phase screening and $5.9\times10^4$ with local fields. The independent-particle rate is too high by 2.3 for a heavy mediator and 3.7 for a light one, and at the large momentum transfers that carry the heavy-mediator rate the local fields, not the head screening, do most of the work. The $f$-sum and the static step bound the rate of any target, loosely for silicon at MeV masses. A mediator that couples to electrons alone breaks the bound on the static step once the lattice relaxes, through the acoustic peak of the electrons that follow the nuclei.

The electron gas at the valence density reproduces the silicon rate within a factor 1.5 between 4.2 and \SI{9.5}{eV}, and fails below the absorption peaks, where it has no gap, and at large energies, where bound electrons carry momenta the gas lacks (\cref{sec:10.3}). Its plasmon is the pole of the inverse dielectric function, but halo dark matter cannot reach it: the slowest plasmon needs a particle faster than \SI{2860}{km.s^{-1}}. A plasmon made otherwise decays into pairs, and into pairs with phonons with a branching estimated at $10^{-3}$ to $10^{-2}$.

Couplings to the electron spin and current need responses beyond the density (\cref{sec:10.4}). Continuity, spin rotation and the product of inversion and time reversal reduce them to twelve real functions per momentum and energy in silicon and germanium, in a singlet block, screened, and a triplet block, unscreened. A velocity coupling adds the transverse current to the density, and a formula that keeps only the dark matter's velocity underestimates its rate in silicon at \SI{10}{MeV} by at least a factor 16 for independent particles, and about 29 with the density row screened.

The electron--hole attraction moves weight toward the gap and adds bound lines below it (\cref{sec:10.5}). In silicon the $1s$ line lies at \SI{1.155}{eV}, below the low-temperature gap of \SI{1.17}{eV}. In a model at \SI{0.91}{keV} the attraction raises the screened loss function by 27 to \SI{29}{\percent} near the edge. An evaluation that drops the antiresonant poles overstates the loss function by a factor 2.12 at small momentum transfer, and dropping the coupling block costs about \SI{10}{\percent}. The triplet channel has the singlet's lines without its local fields or screening.

Core shells enter at large momentum transfer (\cref{sec:10.6}). The germanium $3d$ shell opens at \SI{11.3}{keV} and overtakes the valence at \SI{12.5}{keV}; whether it dominates a rate depends on its final states, and with crystal states it does so above about \SI{30}{MeV} at a one-pair threshold. A filled, flat shell responds as a sum of independent atoms, without interference with the valence. Relativistic corrections to the response are about $10^{-3}$ and at most \SI{2}{\percent} for the shells halo dark matter reaches, and they shift the $f$-sum by $-0.24$ and $-0.76$ percent in silicon and germanium.

Absorbed bosons sample the response at zero momentum (\cref{sec:10.7}). A dark photon is absorbed with the loss function at its mass, suppressed 761 times at \SI{3}{eV} in silicon and enhanced 20 times at the plasmon, \SI{16.4}{eV}. An axion-like particle couples to the spin current, whose response without spin--orbit coupling is three quarters of the current response and is not screened; in a static magnetic field it converts like a dark photon. A scalar couples to the kinetic energy of the electrons, and mixes with the photon only at order $p_\phi^2$ in crystals without a polar axis.

Two insertions of the dark field reach rates that one insertion forbids (\cref{sec:10.8}). The absorption of a particle and its antiparticle is single absorption at twice the mass, and the pseudoscalar pair is absorbed through the axioelectric row without velocity suppression. Dark Compton scattering is the main signal of a dark photon in silicon above \SI{55}{keV}; its seagull form is exact to \SI{1}{\percent} below \SI{165}{keV}, and the published axion-like rate is too large by $(m_\phi + 2m_e)/(m_\phi + m_e)$, 1.836 at \SI{100}{keV}.

Inelastic, exothermic, boosted and fermionic dark matter use the same responses at other kinematics (\cref{sec:10.9}). An endothermic splitting of \SI{1}{eV} raises the lightest mass that makes pairs to \SI{0.60}{MeV}, and an exothermic one turns the spectrum into a halo-independent shell. For boosted particles with a massless mediator the exact density row is 0.65 of the published one at $v = 0.9$, and above $v = \varepsilon_1^{-1/2}$ the transverse row adds Cherenkov emission. An excited state below the gap decays radiatively 41 to 44 times faster inside silicon than in vacuum.

Each approximation of the chapter is a nuisance with a size (\cref{sec:10.10}). At \SI{3}{MeV} with a heavy mediator the silicon rate at the low-temperature gap is $1.85\times10^4$ events per kilogram-year, uncertain by a factor of 1.48 at one standard deviation. The electron--hole attraction and the spread between published independent-particle tables dominate; the gap and the local fields at large $q$ add less than ten percent each. Outside the density rows, the cross term of velocity couplings and the final states of the germanium $3d$ shell are known only between bounds. Throughout the chapter the nuclei were held fixed. \Cref{chap:11} releases them, and the energy a light particle deposits goes into phonons.

